% Options for packages loaded elsewhere
\PassOptionsToPackage{unicode}{hyperref}
\PassOptionsToPackage{hyphens}{url}
%
\documentclass[
]{book}
\usepackage{amsmath,amssymb}
\usepackage{lmodern}
\usepackage{iftex}
\ifPDFTeX
  \usepackage[T1]{fontenc}
  \usepackage[utf8]{inputenc}
  \usepackage{textcomp} % provide euro and other symbols
\else % if luatex or xetex
  \usepackage{unicode-math}
  \defaultfontfeatures{Scale=MatchLowercase}
  \defaultfontfeatures[\rmfamily]{Ligatures=TeX,Scale=1}
\fi
% Use upquote if available, for straight quotes in verbatim environments
\IfFileExists{upquote.sty}{\usepackage{upquote}}{}
\IfFileExists{microtype.sty}{% use microtype if available
  \usepackage[]{microtype}
  \UseMicrotypeSet[protrusion]{basicmath} % disable protrusion for tt fonts
}{}
\makeatletter
\@ifundefined{KOMAClassName}{% if non-KOMA class
  \IfFileExists{parskip.sty}{%
    \usepackage{parskip}
  }{% else
    \setlength{\parindent}{0pt}
    \setlength{\parskip}{6pt plus 2pt minus 1pt}}
}{% if KOMA class
  \KOMAoptions{parskip=half}}
\makeatother
\usepackage{xcolor}
\usepackage{longtable,booktabs,array}
\usepackage{calc} % for calculating minipage widths
% Correct order of tables after \paragraph or \subparagraph
\usepackage{etoolbox}
\makeatletter
\patchcmd\longtable{\par}{\if@noskipsec\mbox{}\fi\par}{}{}
\makeatother
% Allow footnotes in longtable head/foot
\IfFileExists{footnotehyper.sty}{\usepackage{footnotehyper}}{\usepackage{footnote}}
\makesavenoteenv{longtable}
\usepackage{graphicx}
\makeatletter
\def\maxwidth{\ifdim\Gin@nat@width>\linewidth\linewidth\else\Gin@nat@width\fi}
\def\maxheight{\ifdim\Gin@nat@height>\textheight\textheight\else\Gin@nat@height\fi}
\makeatother
% Scale images if necessary, so that they will not overflow the page
% margins by default, and it is still possible to overwrite the defaults
% using explicit options in \includegraphics[width, height, ...]{}
\setkeys{Gin}{width=\maxwidth,height=\maxheight,keepaspectratio}
% Set default figure placement to htbp
\makeatletter
\def\fps@figure{htbp}
\makeatother
\setlength{\emergencystretch}{3em} % prevent overfull lines
\providecommand{\tightlist}{%
  \setlength{\itemsep}{0pt}\setlength{\parskip}{0pt}}
\setcounter{secnumdepth}{5}
\usepackage{booktabs}
\usepackage{amsthm}
\makeatletter
\def\thm@space@setup{%
  \thm@preskip=8pt plus 2pt minus 4pt
  \thm@postskip=\thm@preskip
}
\makeatother

\usepackage{tcolorbox}


\newtcolorbox{blackbox}{
  colback=black,
  coltext=white,
  colframe=black,
  boxsep=5pt,
  arc=4pt}
\newtcolorbox{bonus}{
  colback=blue!15,
  colframe=blue!15,
  coltext=black!80,
  boxsep=5pt,
  arc=4pt}
\newtcolorbox{reflect}{
  colback=green!5,
  colframe=green!5,
  coltext=black!80,
  boxsep=5pt,
  arc=4pt}
\newtcolorbox{assessment}{
  colback=blue!5,
  colframe=blue!5,
  coltext=black!80,
  boxsep=5pt,
  arc=4pt}
\newtcolorbox{progress}{
  colback=purple!10,
  colframe=purple!10,
  coltext=black!80,
  boxsep=5pt,
  arc=4pt}
\newtcolorbox{video}{
  colback=yellow!5,
  colframe=yellow!5,
  coltext=black!80,
  boxsep=5pt,
  arc=4pt}
\newtcolorbox{caution}{
  colback=red!5,
  colframe=red!5,
  coltext=black!80,
  boxsep=5pt,
  arc=4pt}
\newtcolorbox{feedback}{
  colback=black!5,
  colframe=black!5,
  coltext=black!80,
  boxsep=5pt,
  arc=4pt}
\ifLuaTeX
  \usepackage{selnolig}  % disable illegal ligatures
\fi
\usepackage[]{natbib}
\bibliographystyle{apalike}
\IfFileExists{bookmark.sty}{\usepackage{bookmark}}{\usepackage{hyperref}}
\IfFileExists{xurl.sty}{\usepackage{xurl}}{} % add URL line breaks if available
\urlstyle{same} % disable monospaced font for URLs
\hypersetup{
  pdftitle={Philosophy for Life},
  pdfauthor={Andrew Brigham},
  hidelinks,
  pdfcreator={LaTeX via pandoc}}

\title{Philosophy for Life}
\author{Andrew Brigham}
\date{2023-05-13}

\begin{document}
\maketitle

{
\setcounter{tocdepth}{1}
\tableofcontents
}
\hypertarget{course-description}{%
\chapter*{Course Description}\label{course-description}}
\addcontentsline{toc}{chapter}{Course Description}

\hypertarget{course-description-1}{%
\section*{Course Description}\label{course-description-1}}
\addcontentsline{toc}{section}{Course Description}

This course explores philosophy as a way of life accessible to all, in order to think more truthfully, act more justly, and live more faithfully. This course focuses on critical thinking as an invaluable ethical tool for interpreting current events. Students will learn to analyze and evaluate the claims of contemporary culture and religious faith.

\hypertarget{course-navigation}{%
\section*{Course Navigation}\label{course-navigation}}
\addcontentsline{toc}{section}{Course Navigation}

This course is organized into 3 units. Each unit of the course will provide you with the following information:

\begin{itemize}
\tightlist
\item
  A general overview of the topic that will be addressed during the unit.\\
\item
  Specific learning outcomes and topics for the unit.\\
\item
  Resources that you will need to address learning activities and assessments in the unit.\\
\item
  Learning activities both formative and summative will be included in each unit.\\
\item
  Assessments pages will provide you with information about assignments that you will need to complete each for each unit.
\end{itemize}

If you have any questions, do not hesitate to ask. We are here to help and be your guide on this journey.

\hypertarget{course-learning-outcomes}{%
\section*{Course Learning Outcomes}\label{course-learning-outcomes}}
\addcontentsline{toc}{section}{Course Learning Outcomes}

This course will enhance leadership skills and facilitate learning whereby a student will:

\begin{enumerate}
\def\labelenumi{\arabic{enumi}.}
\tightlist
\item
  Develop skills for living a life of wisdom and justice\\
\item
  Cultivate skills of logical argumentation.\\
\item
  Develop and enhance skills of written communication.\\
\item
  Practice the skills of wisdom and reason, such as humility, reflection, charity, rhetoric, and argumentation.\\
\item
  Write a journal that expresses the learned skills concisely and effectively.\\
\item
  Enhance skills of wisdom and reason to live life within the context of justice and faith.\\
\item
  Understand how to become a well-rounded person and citizen.
\end{enumerate}

\hypertarget{texts-and-reading-resources}{%
\section*{Texts and Reading Resources}\label{texts-and-reading-resources}}
\addcontentsline{toc}{section}{Texts and Reading Resources}

\textbf{Required}
- A variety of resources will be provided throughout each unit.

\hypertarget{course-evaluation}{%
\subsection{Course Evaluation}\label{course-evaluation}}

\begin{longtable}[]{@{}
  >{\raggedright\arraybackslash}p{(\columnwidth - 6\tabcolsep) * \real{0.2973}}
  >{\raggedright\arraybackslash}p{(\columnwidth - 6\tabcolsep) * \real{0.2432}}
  >{\raggedright\arraybackslash}p{(\columnwidth - 6\tabcolsep) * \real{0.1892}}
  >{\raggedright\arraybackslash}p{(\columnwidth - 6\tabcolsep) * \real{0.2703}}@{}}
\toprule()
\begin{minipage}[b]{\linewidth}\raggedright
Assignment
\end{minipage} & \begin{minipage}[b]{\linewidth}\raggedright
Activity
\end{minipage} & \begin{minipage}[b]{\linewidth}\raggedright
Value
\end{minipage} & \begin{minipage}[b]{\linewidth}\raggedright
Due
\end{minipage} \\
\midrule()
\endhead
Engagement & Camera, Class, Café & 10\% & Unit 1, 2, 3 \\
Life Skills Video & Personal Reflection Video & 30\% & Unit 1 \\
Arguments Quiz & Quiz & 30\% & Unit 2 \\
Reflection Paper & 750 word Reflection Essay/Reflective Questions & 30\% & Unit 3 \\
\bottomrule()
\end{longtable}

\begin{quote}
For in-person students, check with your teacher about deadlines for assignments.
\end{quote}

Please review the course syllabus to review full list of Course Activities/Requirements/Evaluation.

\emph{Note:} \textbf{WC} = For students required to attend the Writing Centre, you must have a 1:1 writing coach session before you submit your assignment. You must include the writing coach report as page 2 when you submit your assignment.

\hypertarget{academic-integrity-and-avoiding-plagiarism-at-twu}{%
\section*{Academic Integrity and Avoiding Plagiarism at TWU}\label{academic-integrity-and-avoiding-plagiarism-at-twu}}
\addcontentsline{toc}{section}{Academic Integrity and Avoiding Plagiarism at TWU}

As scholars pursuing higher education, academic integrity is a core value of the entire TWU community. Students are invited into this scholarly culture and required to abide by the principles of sound academic scholarship at TWU. This includes, but is not limited to, avoiding all forms of plagiarism and cheating in scholarly work. TWU has a strict policy on plagiarism (see {[}policy below). Learning what constitutes plagiarism and avoiding it is the student's responsibility. Two resources describing plagiarism and how to avoid it have been prepared by TWU Librarian William Badke and are freely available for viewing \href{https://vimeo.com/163566887/bece097c99}{here}.

\begin{quote}
It will be assumed that you have read, understand, and agree to the information provided at the `Academic Dishonesty Policy' button below. If you have any questions at all please contact your instructor.
\end{quote}

{[}button url=``\url{https://www.twu.ca/student-handbook/university-policies/academic-misconduct/procedures-dealing-acts-academic-0}'' target=``\_blank'' label=``Academic Dishonesty Policy'' type=``danger'' classes=``external-link''{]}\_

\url{https://www.twu.ca/about-us/policies-guidelines/university-policies/academic-misconduct-fraud}

\hypertarget{late-assignments-policy}{%
\subsection{Late Assignments Policy}\label{late-assignments-policy}}

Assignments should be submitted on the due date in order to receive full credit. The penalty for unexcused late assignments will be determined by the following scale:

\begin{enumerate}
\def\labelenumi{\arabic{enumi}.}
\tightlist
\item
  For each day or part thereof late, the instructor will reduce the assigned grade by one- third of a letter; e.g., ``A'' to ``A-,'' ``B+'' to ``B.''
\item
  Requests for extensions \textbf{must be communicated in advance} of the due date. Special circumstances may warrant additional consideration.
\end{enumerate}

\hypertarget{writing-centre-sessions}{%
\subsection{Writing Centre Sessions}\label{writing-centre-sessions}}

Please note that you may be required to use the support of our writing center coaches for this course. Plan to book your appointments well in advance, so that the coaches have time to work with you on each of your assignments. The Writing Centre is available to assist all students with their academic writing assignments in any subject at any stage of the writing process. This is a free service. Online Writing Sessions are available, for more information visit \href{http://www.create.twu.ca/learningcommons}{www.create.twu.ca/learningcommons} or contact \href{mailto:writingcentre@twu.ca}{\nolinkurl{writingcentre@twu.ca}}

\hypertarget{program-grade-submission-policy}{%
\subsection{Program Grade Submission Policy}\label{program-grade-submission-policy}}

ALL course work must be submitted for final grading within 14 days following the course end date. The instructor has to submit grades to the Registrar within 14 days of the course end date. In extenuating circumstances, a provision is made for an INC to be entered pending final assignment of a course mark.

\begin{caution}
\emph{Note:} \textbf{WC} = For students required to attend the Writing
Centre.
\end{caution}

\hypertarget{online-learning-orientation}{%
\subsection{Online Learning Orientation}\label{online-learning-orientation}}

For information on how to navigate through this course on Moodle, see the \href{http://create.twu.ca/help/moodle}{\textbf{MOODLE GUIDE}}. Also see the \href{https://www.twu.ca/online-learning}{\textbf{STUDENT GUIDE TO ONLINE LEARNING}} for some tips on how to be successful in the online learning environment.

\hypertarget{course-navigation-1}{%
\section*{Course Navigation}\label{course-navigation-1}}
\addcontentsline{toc}{section}{Course Navigation}

\hypertarget{course-units}{%
\subsection*{Course Units}\label{course-units}}
\addcontentsline{toc}{subsection}{Course Units}

This course is organized into 10 units. Each unit of the course will provide you with the following information:

\begin{itemize}
\tightlist
\item
  A general overview of the key concepts that will be addressed during the unit.
\item
  Specific learning outcomes and topics for the unit.
\item
  Learning activities to help you engage with the concepts. These often include key readings, videos, and reflective prompts.
\item
  The Assessment section provides details on assignments you will need to complete throughout the course to demonstrate your understanding of the course learning outcomes.
\end{itemize}

\begin{caution}
 Note that assessments, including assignments and discussion posts will
 be submitted in Moodle. See the Assessment tab in Moodle for the
 assignment dropboxes.
 \end{caution}

\hypertarget{course-activities}{%
\subsection*{Course Activities}\label{course-activities}}
\addcontentsline{toc}{subsection}{Course Activities}

Below is some key information on features you will see throughout the course.~

\begin{reflect}
 \textbf{\emph{Learning Activity}}\\
 This box will prompt you to engage in course concepts, often by viewing
 resources and reflecting on your experience and/or learning. Most
 learning activities are ungraded and are designed to help prepare you
 for the assessment in this course.
 \end{reflect}

\begin{assessment}
 \textbf{\emph{Assessment}}\\
 This box will signify an assignment or discussion post you will submit
 in Moodle. Note that these demonstrate your understanding of the course
 learning outcomes. Be sure to review the grading rubrics for each
 assignment.
 \end{assessment}

\begin{progress}
 \textbf{\emph{Checking Your Learning}}\\
 This box is for checking your understanding, to make sure you are ready
 for what follows.
 \end{progress}

\begin{video}
 \textbf{\emph{Media}}\\
 This box is for displaying/linking to media, such as videos or songs, in
 order to help illustrate or communicate concepts.
 \end{video}

\begin{caution}
 \textbf{\emph{Note}}\\
 This box signifies key notes, such as where to submit assignments. It
 may also warn you of possible problems or pitfalls you may encounter!
 \end{caution}

\begin{bonus}
 \textbf{\emph{Note}}\\
 This box signifies \ldots another box! Instructors, feel free to add
 your own activity types, such as highlighting case studies, connections
 between topics/learners/instructors, etc.
 \end{bonus}

\begin{feedback}
 \textbf{\emph{Note}}\\
 This box signifies Tips for Instructors. Please delete these before you
 share this course book with your students!
 \end{feedback}

\hypertarget{how-to-navigate-this-book}{%
\subsection*{How To Navigate This Book}\label{how-to-navigate-this-book}}
\addcontentsline{toc}{subsection}{How To Navigate This Book}

To move quickly to different portions of the book, click on the appropriate chapter or section in the table of contents on the left. The buttons at the top of the page allow you to show/hide the table of contents, search the book, change font settings, download a pdf or ebook copy of this book, or get hints on various sections of the book.
\includegraphics{assets/course-intro/menu.png}

The faint left and right arrows at the sides of each page (or bottom of the page if it's narrow enough) allow you to step to the next/previous section. Here's what they look like:
\includegraphics{assets/course-intro/left_arrow.png} \includegraphics{assets/course-intro/right_arrow.png}

\begin{center}\rule{0.5\linewidth}{0.5pt}\end{center}

\hypertarget{writing-standards}{%
\section*{Writing Standards}\label{writing-standards}}
\addcontentsline{toc}{section}{Writing Standards}

For this course, you are expected to follow the writing standards according to APA 7. Please consult the \href{https://owl.purdue.edu/owl/research_and_citation/apa_style/apa_style_introduction.html}{OWL Purdue website} for guidance and seek assistance from the TWU Writing Center and writing coaches as needed. Assignments have rubrics that attribute some marks to APA formatting and cannot be graded as fully meeting expectations if there are APA errors. That said, your conceptual understanding remains of primary importance. It is your responsibility to ensure polished work to the highest standard of which you are capable. This demands meticulous attention to detail, which will become more `natural' with practice. Please seek any necessary clarification from your instructor.

\begin{caution}
\textbf{\emph{It will be assumed that you have read, understand, and
agree to the information provided at the
\href{https://www.twu.ca/student-policies/university-policies/academic-misconduct}{Academic
Dishonesty Policy website}. If you have any questions at all please
contact your instructor.}}
\end{caution}

\hypertarget{course-communities}{%
\chapter*{Course Communities}\label{course-communities}}
\addcontentsline{toc}{chapter}{Course Communities}

\includegraphics{assets/community/luke-porter-NEqEC7qa9FM-unsplash.jpg}

As you begin this course, how will you build community with your fellow learners?

In this course, we have the following tools available to help foster community in your course, including other students who have previously taken this course. Some of these tools will be prescribed and graded (e.g.~Moodle Discussion Forums), others will be up to you to take advantage of.

Check with your course syllabus for which community tools will be used, and consider building your own Community of Practice with your classmates and external colleagues.

\hypertarget{communication-tools}{%
\section*{Communication Tools}\label{communication-tools}}
\addcontentsline{toc}{section}{Communication Tools}

\textbf{Moodle Discussion Forums}: In this course, we ask you to discuss ideas with your colleagues, challenging one another and analyzing key course resources. Refer to the course syllabus for assessment details, as well as the unit Assessment section for discussion questions. Submit your responses in Moodle.

\textbf{Video Conferencing}: We will have scheduled online meetings (Zoom or Teams). Take advantage of these face-to-face conferences! Come prepared with your questions and assigned activities. Refer to the course syllabus and unit activity instructions for details.

\emph{Optional:}

Your cohort may want to engage in other informal discussions to build community and support each other. Consider using the following:

\textbf{Learning Cafe:} This discussion forum in Moodle is a place for you to interact about things going on, share resources, and generally get to know one another. Your posts don't have to be course related. Take this opportunity to connect with fellow learners and learn from one another!

\textbf{Teams:} Every TWU course has a Teams channel, mostly to manage videos. Feel free to use the messaging feature to connect with peers.

\textbf{Twitter hashtag \texttt{\#CRSE\#\#\#}:} You can tweet about this course using \texttt{\#CRSE\#\#\#}.

\textbf{What's App:} Feel free to use a platform that works for you!! What's App is a popular chat forum that learners use for discussions, class projects, etc.

A key takeaway\ldots make these forums work for you! Interact with your peers, learn from each other, and make connections that will stay with you beyond this course.

With that, let's begin the journey together!

\hypertarget{wisdom}{%
\chapter{Wisdom}\label{wisdom}}

\begin{figure}
\centering
\includegraphics{assets/u1/Unit1overview.jpg}
\caption{``Rough waters''}
\end{figure}

CCO Creative Commons by \href{https://isorepublic.com/photo/sun-reflecting-on-waves/}{ISO Republic}

\hypertarget{overview}{%
\section{Overview}\label{overview}}

This unit focuses on wisdom. In this unit, we will (i) introduce some concepts of wisdom found in the literature; and (ii) introduce the ``Life Skills'' to journal and practice. Note that the goal is not to provide a comprehensive summary of each reading, but rather to gain some advice about wisdom from each reading and then practice and journal that advice.

\hypertarget{topics}{%
\section{Topics}\label{topics}}

This unit is divided into the following topics:
1. Wisdom and Action
2. Wisdom and Knowledge
3. Wisdom and Humility
4. Wisdom and Bias
6. Wisdom and Charity
7. Wisdom and Rhetoric

\hypertarget{learning-outcomes}{%
\section{Learning Outcomes}\label{learning-outcomes}}

When you have completed this unit, you should be able to:
- Develop the skills to live a wise and just life
- Develop wisdom, reason, humility, and charity by practicing these skills
- Enhance skills of wisdom and reason to live life within the context of justice and faith.
- Understand how to become a well-rounded person and citizen.

\hypertarget{resources}{%
\section{Resources}\label{resources}}

Nozick, R., 1989, ``What is Wisdom and Why Do Philosophers Love it So?'' in The Examined Life, New York: Touchstone Press, Chapter 23.

\href{https://www.earlymoderntexts.com/assets/pdfs/descartes1644part1.pdf}{Descartes, R., Miller, V. R., \& Miller, R. P. (1983). Principles of philosophy. Dordrecht, Holland: Reidel.}

//todo \#2
\href{https://twu.idm.oclc.org/login?url=https://search.ebscohost.com/login.aspx?direct=true\&db=nlebk\&AN=1085368\&site=eds-live\&scope=site\&ebv=EB\&ppid=pp_1}{Plato, \& Jowett, B. (1999). Apology. Project Gutenberg.}

Stanovish, K. (2020). The Bias that divides us. Quillette. \url{https://quillette.com/2020/09/26/the-bias-that-divides-us/}. (Free subscription)

Johnson, Ralph. 1981. ``Charity Begins at Home,'' Informal Logic Newsletter 3: 4-9.

Heinrichs, J., 2007 ``Set Your Goals,'' chapter 2 in ``Thank you for Arguing: What Aristotle, Lincoln, and Homer Simpson can teach us about the Art of Persuasion. New York : Three Rivers Press,''

\hypertarget{activity-reading}{%
\subsection*{Activity: Reading}\label{activity-reading}}
\addcontentsline{toc}{subsection}{Activity: Reading}

\begin{reflect}
\begin{itemize}
\tightlist
\item
  Robert Nozick, ``What Is Wisdom And Why Do Philosophers Love It So?''\_
  \end{reflect}
\end{itemize}

This topic examines wisdom and action. According to Nozick, wisdom is practical in that wisdom is required to live life well, to cope with and respond to the many challenges of life, and, if possible, avoid some of those challenges. For living life well, wisdom requires both affirming the right beliefs and practicing the right actions. \textbf{How do we know which beliefs and actions are appropriate for living life well?} Nozick argues that a wise person holds beliefs about many aspects of life, including beliefs about life's most important values, beliefs about achieving one's goals, and following the necessary steps to achieve these goals. Likewise, the right actions for living a fulfilling life are derived from applying the aforementioned beliefs. If a person knows about the most important values and goals in life, but never applies them to their behavior, then they are not considered wise. In other words, having knowledge does not entail wisdom. Nozick is also aware that applying the appropriate beliefs and actions is not sufficient to live life well. Sometimes bad things occur for no apparent reason, which may undermine our efforts. A wise person is aware of this problem, and while pursuing their goals, prepares for failure, learns how to avoid it, and how to respond to it if failure occurs.

\hypertarget{highlights-from-the-reading}{%
\subsection{Highlights from the reading:}\label{highlights-from-the-reading}}

\begin{enumerate}
\def\labelenumi{\arabic{enumi}.}
\tightlist
\item
  ``What is wisdom? Wisdom is an understanding of what is important, where this understanding informs a (wise) person's thought and action. Things of lesser importance are kept in proper perspective. (p.~267)
\item
  ``Wisdom is practical; it helps. Wisdom is what you need to understand in order to live well and cope with the central problems and avoid the dangers in the predicaments human beings find themselves in.'' (p.~267)
\item
  ``Wisdom is not just one type of knowledge, but diverse. What a wise person needs to know and understand constitutes a varied list: the most important goals and values of life - the ultimate goal, if there is one; what means will reach these goals without too great a cost; what kinds of dangers threaten the achieving of these goals; how to recognize and avoid or minimize these dangers\ldots{}''(p.~269)
\item
  ``There also will be bits of negative wisdom: certain things are NOT important.'' (p.~269)
  ``A wise person knows these diverse things and lives them. Someone who only knew them, who offered good advice to others yet who lived foolishly himself, would not be termed wise.'' (p.~270)
\item
  ``Wisdom does guarantee success in achieving life's important goals\ldots{}''(p.~270)
\item
  ``A wise person will have gone in the right direction, and if the world thwarts his journey, he will have known how to respond to that too.'' (p.~270)
\item
  ``Moreover, the process of living wisely, pursuing or opening oneself to what is important, taking account of a range of circumstances and utilizing one's fullest capacities to steer skillfully through them, is itself a way of being deeply connected to reality. (p.~276)
  ``The person who lives wisely connects to reality more thoroughly than someone who moves through life spoon-fed by circumstances, even if what they try to feed is reality.'' (p.~276)
\item
  ``Wisdom is not simply knowing how to steer one's way through life, cope with difficulties etc.. It also is knowing the deepest story, being able to see and appreciate the deepest significance of whatever occurs\ldots(p.~276)
\end{enumerate}

\hypertarget{activity-watch-and-reflect}{%
\subsection*{Activity: Watch and Reflect}\label{activity-watch-and-reflect}}
\addcontentsline{toc}{subsection}{Activity: Watch and Reflect}

\begin{reflect}
Please take a moment to watch the following video, a look at the definition of wisdom from philosophers Valerie Tiberius and Philip Kitcher, as well as psychologist Lisa Feldman Barret. The video duration is 5:55 minutes and is an optional resource.
\end{reflect}

This topic contains a learning activity, its purpose is to examine one application of wisdom and setting a personal goal that may improve one's life. Along with setting a personal goal, one must outline the good habits necessary to achieve that goal, and devise strategies to avoid the bad habits that would undermine it. Additionally, one needs to prepare for failure in terms of developing the right habits and reaching the end goal, as well as developing a strategy for avoiding failure and addressing it, if it happens.

\hypertarget{learning-activity-take-notes-for-your-life-journal}{%
\subsection{Learning Activity : Take notes for your Life journal}\label{learning-activity-take-notes-for-your-life-journal}}

Answering the following questions is optional. Answering these questions and saving them as a document will allow you to take notes on the readings in preparation for the Life Journal assignment due at the end of Unit 1.

{[}h5p id=``410''{]}

\hypertarget{activity-reading-1}{%
\subsection*{Activity: Reading}\label{activity-reading-1}}
\addcontentsline{toc}{subsection}{Activity: Reading}

\emph{Pre-reading: (Preface to Principles of Philosophy; Principles 1, 2, 3, 75)}

In the previous topic, we established that knowledge is not sufficient for wisdom because wisdom requires action. That is, knowledge must be applied correctly in order to live life well. While Descartes endorses that view of the limitations of knowledge and the importance of application, he maintains that one should pursue knowledge in the highest domains, such as knowledge of God, science, and morality, because that knowledge benefits the one pursuing wisdom and a fulfilling life.

There are two ideas associated with wisdom and knowledge to discuss. The first idea is that \textbf{acquiring knowledge requires the practice of doubting}. Do not believe everything you read and hear. The second idea is that a \textbf{wise person obtains knowledge of the most valuable things in life, such as knowledge about God, science, and morality}. According to Descartes, pursuing knowledge of these higher domains is not only a good in itself, but that higher knowledge can enable one to live a better life.

Having doubt about something does not entail rejecting its truth. For our purposes, doubting something means applying careful consideration, such as asking yourself: Why should I think this is true? What is the best way to verify the claim? Who can give me an informed opinion? Are there any concerns about this content? Which aspects of the content are strong and which are weak? Healthy doubt is valuable for acquiring knowledge and living life well. For example, whenever we detect incoming scams, such as phone and computer scams, we employ doubt. In this regard, one intends to uncover the nature of the scam and warn others of its harmful potential. The hesitation (or doubt) to believe and act according to the computer scam is similar to the hesitation one should apply when listening to the news, a podcast, a sermon, or reading an opinion on Facebook or Twitter. However, particularly in ordinary affairs, Decartes also warns that extreme doubt may be harmful. Neither the need for basic necessities nor the importance of friendship should be questioned. When individuals are faced with situations such as these, it is often necessary that they act even if they lack all the information about the situation. Descartes argues that remaining paralyzed by doubt is an inadvisable course.

To pursue a higher knowledge about God, science, and morality, Descartes imagines the structure of a tree, whose roots are metaphysics, the trunk is physics, and the branches represent the remainder of scientific inquiry, including morality (note that `science' in this context encompasses more domains than a contemporary definition of science). Descartes believed that in order to understand the human condition, one must have an understanding of metaphysics and science. Descartes' point about seeking this type of valuable knowledge about God, science, and morality is relevant because of the human obsession with entertainment. We often seek entertainment and pleasure rather than searching for the truth about the most important ideas in life. That obsession with entertainment and the neglect of seeking valuable knowledge undermines one's potential to gain wisdom and seek a fulfilling life.

\hypertarget{highlights-from-the-reading-preface}{%
\subsection{Highlights from the reading (Preface):}\label{highlights-from-the-reading-preface}}

\begin{enumerate}
\def\labelenumi{\arabic{enumi}.}
\tightlist
\item
  For example, the word `philosophy' means the study of wisdom, and by `wisdom' it means not only prudence in our everyday affairs but also a perfect knowledge of all things that mankind is capable of knowing\ldots'' (179)
\item
  ``\ldots we should not doubt the probable truths which concern the conduct of life\ldots{}''(182).
\item
  ``Thus, the whole of philosophy is like a tree. The roots are metaphysics, the trunk is physics, and the branches emerging from the trunk are all the other sciences\ldots{}'' (186)
\item
  ``\ldots For although the truth often does not touch our imagination as much as falsehood and pretense, because it seems less striking and more plain, nevertheless the satisfaction it produces is always more durable and more solid.'' (188)
\item
  ``The second benefit {[}of studying Descartes' principles of philosophy{]} is that the study of these principles will accustom people little by little to form better judgments about all the things they come across, and hence will make them wiser.'' (188)
\end{enumerate}

Please take a moment to view this optional video. In this video, the Honorable Richard Holloway, who is a writer, broadcaster, and cleric from Scotland, seeks to initiate a debate around the more profound spiritual questions of life. The duration of this video is 4:45 minutes.
\href{https://youtu.be/LERi_Xjfio4}{plugin:youtube}

The aim of the following learning activities is to explore and learn about these two applications of knowledge: the practice of doubt and the acquisition of valuable knowledge over entertainment.

\begin{center}\rule{0.5\linewidth}{0.5pt}\end{center}

\hypertarget{learning-activity-1.2-learning-to-doubt---simulation-exercise-1}{%
\subsection{Learning Activity 1.2: Learning to doubt - Simulation Exercise 1}\label{learning-activity-1.2-learning-to-doubt---simulation-exercise-1}}

:fa-puzzle-piece: The simulation exercises are meant to simulate a conversation and introduce students to the various topics in Unit 1. Students are encouraged to not only review the simulation exercises but to practice the applications of wisdom in their daily lives.

{[}h5p id=``409''{]}

\hypertarget{learning-activity-1.3-take-notes-for-your-life-journal}{%
\subsection{Learning Activity 1.3: Take notes for your Life Journal}\label{learning-activity-1.3-take-notes-for-your-life-journal}}

:fa-pencil: Answering the following questions is optional. Answering these questions and saving them as a document will allow you to take notes on the readings in preparation for the Life Journal assignment due at the end of Unit 1.

{[}h5p id=``412''{]}

\emph{Pre-reading: \url{http://classics.mit.edu/Plato/apology.html}}

The third topic focuses on wisdom and humility. Our previous two topics discussed the importance of knowledge when seeking wisdom. Nevertheless, even Descartes recognizes that knowledge in some domains is difficult to obtain. Plato believed that a wise person understands their limitations about what they know and do not know. In other words, \textbf{a wise person is aware of their ignorance}. The word ``ignorance'' in this case does not have a negative connotation, but rather describes a lack of knowledge about a particular thing.

In ``the Apology,'' Plato introduces the character Socrates, who is on trial for his life. During his defense, Socrates describes the account of his friend, Chaerephon, seeking answers from the Oracle of Delphi, wondering if there is any person wiser than Socrates. In reply to Chaerephon, the oracle says that no one is wiser than Socrates. Socrates is puzzled by the revelation about his wisdom since he considers himself to be unwise. Socrates then asks questions of different people in the city, such as politicians, poets, and artisans, in an attempt to solve the mystery about the oracle's revelation. Socrates discovers that while these individuals claim to know a great deal about issues such as the nature of beauty and what is good, he observes that they know very little about such issues. These people think they possess knowledge, when in fact they do not, and that is why they lack wisdom. Their way of thinking contrasts with Socrates, who thinks he lacks knowledge about these issues; more specifically, Socrates is aware of his ignorance about the nature of beauty and what is good, while the others are unaware of their own ignorance about these same issues. Because of this, Socrates is the wisest person.

One application of being aware of ignorance for seeking wisdom is to practice a humble attitude towards the complex issues in life, such as how to interpret the Bible, evolution, the nature of sexuality, the implications of social media for our lives, and was Donald Trump an effective president? These types of questions represent complex ideas that even some of the most specialized scholars cannot answer. A wise person demonstrates their humility towards these complex issues by asking thoughtful questions about the opposing view, knowing when to withhold judgment and opinion, and with a willingness to change their mind.

\hypertarget{highlights-from-the-reading-1}{%
\subsection{Highlights from the reading:}\label{highlights-from-the-reading-1}}

``When I heard the answer, I said to myself, What can the god mean? and what is the interpretation of this riddle? for I know that I have no wisdom, small or great. What can he mean when he says that I am the wisest of men? And yet he is a god and cannot lie; that would be against his nature. After a long consideration, I at last thought of a method of trying the question. I reflected that if I could only find a man wiser than myself, then I might go to the god with a refutation in my hand. I should say to him,''Here is a man who is wiser than I am; but you said that I was the wisest.'' Accordingly I went to one who had the reputation of wisdom, and observed to him - his name I need not mention; he was a politician whom I selected for examination - and the result was as follows: When I began to talk with him, I could not help thinking that he was not really wise, although he was thought wise by many, and wiser still by himself; and I went and tried to explain to him that he thought himself wise, but was not really wise; and the consequence was that he hated me, and his enmity was shared by several who were present and heard me. So I left him, saying to myself, as I went away: Well, although I do not suppose that either of us knows anything really beautiful and good, I am better off than he is - for he knows nothing, and thinks that he knows. I neither know nor think that I know. In this latter particular, then, I seem to have slightly the advantage of him. Then I went to another, who had still higher philosophical pretensions, and my conclusion was exactly the same. I made another enemy of him, and of many others besides him.''

``After this I went to one man after another, being not unconscious of the enmity which I provoked, and I lamented and feared this: but necessity was laid upon me - the word of God, I thought, ought to be considered first. And I said to myself, Go I must to all who appear to know, and find out the meaning of the oracle. And I swear to you, Athenians, by the dog I swear! - for I must tell you the truth - the result of my mission was just this: I found that the men most in repute were all but the most foolish; and that some inferior men were really wiser and better. I will tell you the tale of my wanderings and of the''Herculean'' labors, as I may call them, which I endured only to find at last the oracle irrefutable. When I left the politicians, I went to the poets; tragic, dithyrambic, and all sorts. And there, I said to myself, you will be detected; now you will find out that you are more ignorant than they are. Accordingly, I took them some of the most elaborate passages in their own writings, and asked what was the meaning of them - thinking that they would teach me something. Will you believe me? I am almost ashamed to speak of this, but still I must say that there is hardly a person present who would not have talked better about their poetry than they did themselves. That showed me in an instant that not by wisdom do poets write poetry, but by a sort of genius and inspiration; they are like diviners or soothsayers who also say many fine things, but do not understand the meaning of them. And the poets appeared to me to be much in the same case; and I further observed that upon the strength of their poetry they believed themselves to be the wisest of men in other things in which they were not wise. So I departed, conceiving myself to be superior to them for the same reason that I was superior to the politicians.''

``At last I went to the artisans, for I was conscious that I knew nothing at all, as I may say, and I was sure that they knew many fine things; and in this I was not mistaken, for they did know many things of which I was ignorant, and in this they certainly were wiser than I was. But I observed that even the good artisans fell into the same error as the poets; because they were good workmen they thought that they also knew all sorts of high matters, and this defect in them overshadowed their wisdom - therefore I asked myself on behalf of the oracle, whether I would like to be as I was, neither having their knowledge nor their ignorance, or like them in both; and I made answer to myself and the oracle that I was better off as I was.''

The aim in the following learning activities is to practice humility about what we know and don't know. One way to practice humility is to ask questions and limit how often we share our opinions.

Please take a moment to watch this optional video with Dr.~Ian Church, What is intellectual humility? A doxastic account. Here is the first of three parts. If you have the time, you should watch them all. Please note that this video lasts for approximately 2:27 minutes.

\href{https://youtu.be/8CZIkGEJYRY}{plugin:youtube}

\hypertarget{learning-activity-1.4-practicing-humility---simulation-exercise-2}{%
\subsection{Learning Activity 1.4: Practicing Humility - Simulation Exercise 2}\label{learning-activity-1.4-practicing-humility---simulation-exercise-2}}

:fa-puzzle-piece: The simulation exercises are meant to simulate a conversation and introduce students to the various topics in Unit 1. Students are encouraged to not only review the simulation exercises but to practice the applications of wisdom in their daily lives.

{[}h5p id=``424''{]}

\hypertarget{learning-activity-1.5-take-notes-for-your-life-journal}{%
\subsection{Learning Activity 1.5: Take notes for your Life Journal}\label{learning-activity-1.5-take-notes-for-your-life-journal}}

:fa-pencil: Answering the following questions is optional. Answering these questions and saving them as a document will allow you to take notes on the readings in preparation for the Life Journal assignment due at the end of Unit 1.

{[}h5p id=``427''{]}

\emph{Pre-reading: \url{https://quillette.com/2020/09/26/the-bias-that-divides-us/}}

In the previous topic, we discussed the importance of humility when seeking knowledge. Topic 4 explores the relationship between wisdom and bias and this relationship is valuable for seeking humility. Recall that a wise person seeking humility recognizes their own ignorance regarding the complex issues in life. Identifying personal ignorance also requires acknowledging one's biases. When one recognizes the deep influence of human biases on knowledge acquisition, it becomes easier to recognize one's ignorance.

There are many types of biases, both social and cognitive. \textbf{Social bias} is a kind of negative attitude and opinion toward other members in society, usually without good reason. Examples of social bias include prejudice toward other members of society based on their race, gender, age, and so forth. A \textbf{cognitive bias} refers more generally to human intelligence, and, in some cases, to human reasoning. Specifically, cognitive bias occurs when personal interest conflicts with the acquisition of knowledge; that is, humans tend to make errors of judgment due to their bias towards self-interest. One form of cognitive bias is called \textbf{confirmation bias}. Confirmation bias is when people attempt to acquire knowledge by consulting evidence that only supports their own agenda. An example of this would be a person who believes that God exists and supports their position by viewing sources that affirm the existence of God and ignoring or misinterpreting sources that deny this belief. That person seeks information that reinforces their current opinions and generally avoids information that undermines their current opinions.

A more precise way to think about confirmation bias is called ``myside'' bias. While myside bias includes the error of favoring evidence that supports one's previously held opinions, myside bias locates that error within the context of group membership. Those previously held opinions often reflect convictions to agree with one's social group and discontentment towards agreeing with opposing social groups. For example, Stanovich describes a person who commits what is perceived as a virtuous act. If that person is a member of your social group, you will be more inclined to affirm the act as virtuous. If that person is not a member of your social group, you will be less inclined to affirm the act as virtuous. This error of reasoning comes from affirming the virtue of the act only by association with a particular social group. Stanovich contends that myside bias affects everyone, especially intellectual elites, and those who believe they are immune to such biases are more susceptible to them.

\hypertarget{highlights-from-the-reading-2}{%
\subsection{Highlights from the reading:}\label{highlights-from-the-reading-2}}

``\ldots myside bias: people evaluate evidence, generate evidence, and test hypotheses in a manner biased toward their own prior beliefs, opinions, and attitudes.''

``Research has shown that myside bias is displayed in a variety of experimental situations: people evaluate the same virtuous act more favourably if committed by a member of their own group and evaluate a negative act less unfavourably if committed by a member of their own group; they evaluate an identical experiment more favourably if the results support their prior beliefs than if the results contradict their prior beliefs; and when searching for information, people select information sources that are likely to support their own position. Even the interpretation of a purely numerical display of outcome data is tipped in the direction of the subject's prior belief. Likewise, judgments of logical validity are skewed by people's prior beliefs. Valid syllogisms with the conclusion ``therefore, marijuana should be legal'' are easier for liberals to judge correctly and harder for conservatives; whereas valid syllogisms with the conclusion ``therefore, no one has the right to end the life of a fetus'' are harder for liberals to judge correctly and easier for conservatives''

``myside bias is one of the most ubiquitous of biases because it is exhibited by the vast majority of subjects studied. Myside bias is also not limited to individuals with certain cognitive or demographic characteristics. It is one of the most universal of cognitive biases.''

``In short, the convictions that determine your side when you think in a mysided fashion, often don't come from rational thought.''

``However, one particular bias---myside bias---sets a trap for the cognitively sophisticated. Regarding most biases, they are used to thinking---rightly---that they are less biased. However, myside thinking about your political beliefs represents an outlier bias where this is not true.''

``This may lead to a particularly intense bias blind spot among certain cognitive elites. If you are a person of high intelligence, if you are highly educated, and if you are strongly committed to an ideological viewpoint, you will be highly likely to think you have thought your way to your viewpoint.''

``University faculty in the social sciences fit this bill perfectly. And the opening for a massive bias blind spot occurs when these same faculty think that they can objectively study, within the confines of an ideological monoculture, the characteristics of their ideological opponents.''

``Anything that makes us more skeptical about our beliefs will tend to decrease the myside bias that we display (by preventing beliefs from turning into convictions).''

``In short, just as we are gorging on fat-laden food that is not good for us because our bodies were built by genes with a selfish replicator survival logic, so we are gorging on memes that fit our resident beliefs because cultural replicators have a similar survival logic. And just as our overconsumption of fat-laden fast foods has led to an obesity epidemic, so our overconsumption of congenial memes has made us memetically obese as well.''

``When I talk to lay audiences about different types of cognitive processes, I use the example of broccoli and ice cream. Some cognitive processes are demanding but necessary. They are the broccoli. Other thinking tendencies come naturally to us and they are not cognitively demanding processes. They are the ice cream. In lectures, I point out that broccoli needs a cheerleader, but ice cream does not. This is why education rightly emphasizes the broccoli side of thinking---why it stresses the psychologically demanding types of thinking that people need encouragement to practice.''

When you first wake up, remember that you are biased. As you discuss controversies with fellow students, colleagues, family, and friends, your biases are hard at work, practically bullying you into thinking a certain way. The moment you deny this reality, you succumb to yet another bias. The aim in the following learning activity is to become more aware of one's cognitive biases and how that awareness can promote an attitude of humility.

Watch the optional video below. Jonathan Haidt, a renowned social psychologist, shares a personal example of confirmation bias on social media and emphasizes the importance of listening to ideas from different political perspectives. Please note that the video is 2:13 minutes in length.

\href{https://youtu.be/XO0AVyF6EVg}{plugin:youtube}

\hypertarget{learning-activity-1.6-awareness-of-bias---simulation-exercise-3}{%
\subsection{Learning Activity 1.6: Awareness of bias - Simulation Exercise 3}\label{learning-activity-1.6-awareness-of-bias---simulation-exercise-3}}

:fa-puzzle-piece: The simulation exercises are meant to simulate a conversation and introduce students to the various topics in Unit 1. Students are encouraged to not only review the simulation exercises but to practice the applications of wisdom in their daily lives.

{[}h5p id=``419''{]}

\hypertarget{learning-activity-1.7-take-notes-for-your-life-journal}{%
\subsection{Learning Activity 1.7: Take notes for your Life Journal}\label{learning-activity-1.7-take-notes-for-your-life-journal}}

:fa-pencil: Answering the following questions is optional. Answering these questions and saving them as a document will allow you to take notes on the readings in preparation for the Life Journal assignment due at the end of Unit 1.

{[}h5p id=``429''{]}

\emph{Pre-reading: Ralph H. Johnson, ``Charity Begins at Home''}

In Topic 5, the emphasis is on wisdom and charity. The term ``charity'' in this case refers to \textbf{a principle of treating a person's argument and/or opinion in its strongest form}. In the previous two topics, we discussed humility and bias, both of which are relevant to practicing charity. For example, when employing humility, one is aware of their ignorance and the need to ask questions about the opponent's view. That questioning can be a type of charity when done for the purpose of gaining understanding. Similarly, one is aware of myside bias, and that bias influences the interpretation and treatment of the opponent, especially if that opponent is from another social group. Recognizing the powerful influence of myside bias promotes charity by helping constrain one's immediate rejection of the opponent's view.

Charity involves more than humility and awareness of bias: it requires willpower. One must willfully, and temporarily, embrace the opponent's view, assist with constructing that view in its best possible light, and then express the view better than the opponent. The act of embracing an opposing argument is challenging for many reasons, but, as stated in the previous chapter on bias, humans are not built to act with charity toward those with whom they disagree, especially when those convictions are strong. Put in a different way, humans struggle to, on the one hand, affirm a viewpoint in which they hold strong convictions, and then, on the other hand, compartmentalize that view away while attempting to charitably evaluate objections to that conviction.

In his article, ``Charity Begins at Home,'' Johnson outlines the general principle of charity and then discusses some of the challenges with that principle. Practicing charity involves three steps:

\begin{enumerate}
\def\labelenumi{\arabic{enumi}.}
\tightlist
\item
  Identifying the argument: Does the argument even exist? According to Johnson, the first step is difficult in its own right because the art of constructing a well-organized argument is in decline.
\item
  Reconstructing the argument: This step requires that one re-build the argument to follow the rules of logical inference. Sometimes arguments include additional hidden premises, irrelevant information, and are simply weak or invalid. One must reconstruct the argument in its strongest form.
\item
  Presenting the main objections to the argument: Now that the argument has been identified and properly reconstructed, one should focus on the criticisms that undermine the main premises of the argument and avoid criticizing irrelevant information and implications of the argument.
\end{enumerate}

According to Johnson, a major problem with the principle of charity is that it requires too much effort and time to apply, as few people can construct an argument correctly in the first place. His solution to this problem is to apply the principle only when the following conditions are met: (i) the argument is fully expressed; (ii) the argument is brought forward by a serious person; (iii), the argument is about a serious issue. One implication of those criteria for the principle is that one should apply charity only when there is reason to think that the opponent will benefit from that charity.

\hypertarget{highlights-from-the-reading-3}{%
\subsection{Highlights from the reading:}\label{highlights-from-the-reading-3}}

``The Principle of Charity which governs all levels of argument analysis is that the critic should provide the best possible interpretation of the material under consideration.'' (p.~5)

``The justification of the principle is\ldots ethical. One is under the general obligation to be fair in one's dealings with others\ldots{}'' (p.5)

``There is also a prudential reason for adhering to the principle\ldots{} It tells you that you want to interpret the argument's meaning in whatever way makes the most sense and force out of it, because otherwise, it can easily be reformulated\ldots to meet your objections.'' (p.5)

``The first step in argument analysis must be to decide whether or not the passage under scrutiny contains an argument\ldots{}''(p.5)

``Once a passage has been identified as containing an argument we are then faced with the task of identifying and setting forth its premises and conclusions in their logical relationships.'' (p.5)

``Once the argument has been identified, extracted, and reconstructed, the job of the critic can begin\ldots one should look for the strongest possible criticisms of the argument\ldots{}''(p.6)

``Effective argument analysis can be a grueling task.'' (p.8)

``Many times I have found myself wrestling with diffuse and poorly expressed arguments that seem to have been sloppily put together - something dashed off in five minutes by a loose reasoner with scant knowledge of the demands of the argumentative process.'' (p.8)

``In defense of this proposal {[}the proposal that the only arguments deserving charity are those that meet the three criteria{]}, let me make two points. First, the individual who dashes off a poorly expressed argument is, in all probability, not going to derive any profit from the critics' laborious undertaking. His or her commitment to the rational process is too slight and powers of reasoning too underdeveloped (else we would not be in this situation in the first place). (p.~8).

The aim in the following learning activity is to learn about practicing charity. The goal is to embrace your opponent's position, assist them in making their position as strong as possible, and only then introduce objections to their position.

You are invited to watch the following optional video. John Corvino who is a speaker, writer, philosophy professor, and Dean of the Irvin D. Reid Honors College at Wayne State University in Detroit, is sharing about The Principle of Charity. Please note that this video is 2:07 minutes long.

\href{https://youtu.be/LZZ7tQnI2-M}{plugin:youtube}

\hypertarget{learning-activity-1.8-acting-with-charity---simulation-exercise-4}{%
\subsection{Learning Activity 1.8: Acting with charity - Simulation Exercise 4}\label{learning-activity-1.8-acting-with-charity---simulation-exercise-4}}

:fa-puzzle-piece: The simulation exercises are meant to simulate a conversation and introduce students to the various topics in Unit 1. Students are encouraged to not only review the simulation exercises but to practice the applications of wisdom in their daily lives.

{[}h5p id=``422''{]}

\hypertarget{learning-activity-1.9-take-notes-for-your-life-journal}{%
\subsection{Learning Activity 1.9: Take notes for your Life Journal}\label{learning-activity-1.9-take-notes-for-your-life-journal}}

:fa-pencil: Answering the following questions is optional. Answering these questions and saving them as a document will allow you to take notes on the readings in preparation for the Life Journal assignment due at the end of Unit 1.

{[}h5p id=``444''{]}

\emph{Pre-reading: Jay Heinrichs, ``Set Your Goals,'' chapter 2 in Thank you for Arguing: What Aristotle, Lincoln, and Homer Simpson can teach us about the Art of Persuasion.''}

Topic 6 is about wisdom and rhetoric. Rhetoric refers to \textbf{the art of persuasion} or influence. Previously we discussed how wisdom requires an awareness and application of our actions, knowledge, humility, biases, and charity. The art of persuasion employs all of these skills. Along with those skills of wisdom, rhetoric also requires some tools of argumentation. A clarification regarding rhetoric is that the terms ``argument'', ``argue'', and ``argumentation'' can mean different things in different contexts. Students will learn the technical definitions of ``argument'' and ``argumentation'' in Unit 2, which may differ from how those terms are used in this topic of rhetoric. The student will also have the opportunity to develop some of the tools of technical argumentation in Unit 2, which will hone their rhetorical skills.

Arguments with other people, such as family, friends, or colleagues, are part of everyday life. A wise person knows how to argue without fighting. According to Heinrichs, a necessary skill for arguing productively is knowing the appropriate goal you wish to bring about. Are you trying to influence someone's decision? Is that feasible? Are you hoping to feel like a winner? Is that an appropriate goal? It may be your goal to change the mood of your opponent, or simply to inform your opponent of a different point of view. Would you like to end the argument quickly? After the argument, would you like to maintain the relationship with the person? Should you pursue either of these goals? How will you go about accomplishing your goal? In what way are you prepared to compromise? Are you prepared to experience defeat from your opponent's point of view in order to achieve your goal? What is the purpose of the argument and how should it be achieved?

Exercising the previous skills of wisdom when arguing with other people is helpful for practicing rhetoric. Wise actions, for example, require setting reasonable goals and avoiding the pitfalls that may undermine these goals. Similarly, when applying rhetoric in an argument with someone else, focus on the goal you wish to bring about and be willing to concede as much as possible to get there. Furthermore, remember that wisdom requires knowledge and humility. When those skills are applied to rhetoric, demonstrate the ability to challenge your own views to show your intention to understand the opponent's position. Knowledge of one's own bias is also essential to rhetoric. Ensure that the opponent understands that you are aware of your own biases and that you will listen to and learn from them. And in everything, act with charity. When arguing with others, assist them with their arguments, and, if possible, suggest ways in which their arguments can be strengthened. Demonstrate your comfort with the exchange by temporarily embracing their viewpoint. Charity of this kind encourages your opponent to act in a similar manner, and, in doing so, may reduce the defensiveness of each person. Depending on your rhetorical goal when arguing with another person, the tactics above may be useful.

\hypertarget{highlights-from-the-reading-4}{%
\subsection{Highlights from the reading:}\label{highlights-from-the-reading-4}}

``This chapter will help you distinguish between an argument and a fight, and to choose what you want to get out of an argument.'' (p, 16)
``You succeed in an argument when you persuade your audience.'' (p.~16)

``\ldots argument by the stick - fighting - is no argument. It never persuades, it only inspires revenge or retreat.'' (p.17)

``In a fight, one person takes out his aggression on another. Donald Trump was fighting when he said of Rosie O'Donnell, ``I mean, I'd look at her in that fat, ugly face of hers, I'd say `Rosie, you're fired.' On the other hand, when Geroge Foreman tries to sell you a grill, he makes an argument: persuasion that tries to change your mood, your mind, or your willingness to do something.'' (p.17)

``The basic difference between an argument and a fight: an argument, done skillfully, gets people to want what you want. You fight to win; you argue to achieve agreement.'' (p.~17)

``Changing the mood is the easiest goal, and usually the one you work on first {[}\#1 goal{]}.'' (p.~21)
`` \ldots making them decide what you want.'' {[}\#2 goal{]} (p.~21)

``Goal number three - in which you get an audience to do something or stop doing it - is the most difficult.'' {[}\#3 goal{]} (p.22).

The aim in the following learning activity is to practice a skill of rhetoric: set a goal of changing the mood of the argument and try to outline the way of achieving that goal. One way to change the mood of the argument is by (implicitly) affirming your opponent's view. Say something like this: ``Fair enough. I need to think about this. Can you explain that in more detail?'' OR ``That's good advice. Can you tell me a bit more?'' Be careful not to patronize, but sincerely affirm your opponent's view to change the mood of the debate, from each person being defensive to each person being receptive.

Take a moment to watch this optional video. Here, Camille A. Langston, a specialist in rhetoric and nineteenth-century women, explains the basics of deliberative rhetoric and shares tips for appealing to the ethos, logos, and pathos of your audience. This video is 4:29 minutes long.

\href{https://youtu.be/3klMM9BkW5o}{plugin:youtube}

\hypertarget{learning-activity-1.9-the-art-of-persuasion---simulation-exercise-5}{%
\subsection{Learning Activity 1.9: The art of persuasion - Simulation Exercise 5}\label{learning-activity-1.9-the-art-of-persuasion---simulation-exercise-5}}

:fa-puzzle-piece: The simulation exercises are meant to simulate a conversation and introduce students to the various topics in Unit 1. Students are encouraged to not only review the simulation exercises but to practice the applications of wisdom in their daily lives.

{[}h5p id=``420''{]}

\hypertarget{learning-activity-1.10-take-notes-for-your-life-journal}{%
\subsection{Learning Activity 1.10: Take notes for your Life Journal}\label{learning-activity-1.10-take-notes-for-your-life-journal}}

:fa-pencil: Answering the following questions is optional. Answering these questions and saving them as a document will allow you to take notes on the readings in preparation for the Life Journal assignment due at the end of Unit 1.

{[}h5p id=``428''{]}

\hypertarget{assessment}{%
\section{Assessment}\label{assessment}}

\hypertarget{grading-criteria}{%
\subsubsection{Grading Criteria}\label{grading-criteria}}

\hypertarget{submission-instructions}{%
\subsubsection{Submission Instructions}\label{submission-instructions}}

Submit your Reflective Journal to the Moodle dropbox under Assessments.

\hypertarget{forum-introductions}{%
\subsection{Forum Introductions}\label{forum-introductions}}

Also remember to introduce yourself to your TWU instructor by writing a brief introduction in the online Discussion Forum. Please share about yourself, such as where you live, what you are studying, and the kind of things that interest you. Perhaps some questions you have about this course? Share a picture of someone or something that means a lot to you (i.e.~your pet, family, favourite movie or book, etc.).

\hypertarget{checking-your-learning}{%
\section{Checking Your Learning}\label{checking-your-learning}}

Before you move on to the next unit, check that you are able to:\\
{[}fa=check /{]} \ldots{}\\
{[}fa=check /{]} \ldots{}\\
{[}fa=check /{]} \ldots{}

\hypertarget{title}{%
\chapter{Title}\label{title}}

\hypertarget{title-1}{%
\chapter{Title}\label{title-1}}

\hypertarget{title-2}{%
\chapter{Title}\label{title-2}}

\hypertarget{title-3}{%
\chapter{Title}\label{title-3}}

\hypertarget{title-4}{%
\chapter{Title}\label{title-4}}

\hypertarget{title-5}{%
\chapter{Title}\label{title-5}}

\hypertarget{title-6}{%
\chapter{Title}\label{title-6}}

\hypertarget{assessment-1}{%
\chapter*{Assessment}\label{assessment-1}}
\addcontentsline{toc}{chapter}{Assessment}

The following assignments are opportunities for learners to demonstrate their understanding of the course outcomes. Please confirm assignment details with your instructor, referring to the course syllabus.

Note that Assignment dropboxes are located in Moodle. Also refer to the Course Schedule in Moodle for the specific due dates.

\hypertarget{assignment}{%
\section*{Assignment:}\label{assignment}}
\addcontentsline{toc}{section}{Assignment:}

\begin{assessment}

\end{assessment}

\hypertarget{grading-criteria-1}{%
\subsection*{Grading Criteria}\label{grading-criteria-1}}
\addcontentsline{toc}{subsection}{Grading Criteria}

See the following rubric that explains how your assignment will be evaluated. Also available as a \href{assets/assessment/Identity-as-a-Teacher-RUBRIC.pdf}{pdf}

\#\#\#\# APA/WRITING \{-\}

\textbf{Unsatisfactory:} Paper does not model language and conventions used in scholarly literature. Writing is not well-organized. Several errors in grammar or composition. Sources are not cited. APA citations are not appropriately formatted.

\textbf{Developing:} Paper partially models language and conventions used in scholarly literature. Writing is somewhat well organized and includes some errors in grammar or composition. Not all sources cited. APA citations are generally formatted correctly, with several errors.

\textbf{Proficient:} \emph{Paper consistently models language and conventions used in scholarly literature. Writing is well-organized and includes few (if any) errors in grammar or composition. All resources are appropriately cited (including in-text citations and bibliography information). Few (if any) errors in APA citations.}

\textbf{Exemplary:} Paper is an exemplar of language and conventions used in scholarly literature. Writing is well-organized and free of errors in grammar or composition. All resources are appropriately cited. No errors in APA format.

\hypertarget{statement-of-teaching-identity}{%
\subsubsection*{STATEMENT OF TEACHING IDENTITY}\label{statement-of-teaching-identity}}
\addcontentsline{toc}{subsubsection}{STATEMENT OF TEACHING IDENTITY}

\textbf{Unsatisfactory:} Does not provide a statement about identity as a teacher/facilitator

\textbf{Developing:} Provides an unclear statement about identity as a teacher/facilitator.

\textbf{Proficient:} \emph{Provides a clear, concise, and powerful statement about identity as a teacher/facilitator.}

\textbf{Exemplary:} Provides a clear, concise, and powerful statement about identity as a teacher/facilitator. Statement incorporates theory or research from course materials.

\hypertarget{developing-a-cohesive-and-logical-academic-argument}{%
\subsubsection*{DEVELOPING A COHESIVE AND LOGICAL ACADEMIC ARGUMENT}\label{developing-a-cohesive-and-logical-academic-argument}}
\addcontentsline{toc}{subsubsection}{DEVELOPING A COHESIVE AND LOGICAL ACADEMIC ARGUMENT}

\textbf{Unsatisfactory:} Does not make a focused, cohesive, or logical academic argument. Paper is confusing, and is missing an introduction, body, or conclusion. Transitions between sections and ideas are missing.

\textbf{Developing:} Makes an academic argument that is only partially focused, cohesive and logical. Paper is generally organized, but is missing an introduction, body, or conclusion. Transitions between sections and ideas are unclear.

\textbf{Proficient:} \emph{Makes a focused, cohesive, logical academic argument. Paper is effectively organized and includes an introduction, body, and conclusion. Transitions between sections and ideas are clear.}

\textbf{Exemplary:} Makes a focused, cohesive, logical and compelling academic argument. Paper is effectively organized and includes an introduction, body, and conclusion. Transitions between sections and ideas are clear, and build on each other.

\hypertarget{analysis-of-identity-as-a-teacher}{%
\subsubsection*{ANALYSIS OF IDENTITY AS A TEACHER}\label{analysis-of-identity-as-a-teacher}}
\addcontentsline{toc}{subsubsection}{ANALYSIS OF IDENTITY AS A TEACHER}

\textbf{Unsatisfactory:} Does not include three important aspects of identity as a teacher/facilitator. Does not include an analysis.

\textbf{Developing:} Lists but does not discuss three important aspects of identity as a teacher/facilitator. Includes a partial analysis.

\textbf{Proficient:} \emph{Includes a detailed discussion of three important aspects of identity as a teacher/facilitator. Includes thoughtful analysis of each of the three elements.}

\textbf{Exemplary:} Includes a detailed discussion of three important aspects of identity as a teacher/facilitator. Includes a thoughtful analysis, integrating scholarly literature to support analysis and furthering scholarly thinking related to teacher identity.

\hypertarget{scholarly-integration}{%
\subsubsection*{SCHOLARLY INTEGRATION}\label{scholarly-integration}}
\addcontentsline{toc}{subsubsection}{SCHOLARLY INTEGRATION}

\textbf{Unsatisfactory:} Does not integrate references to support claims and assertions made in the paper.

\textbf{Developing:} Integrates references to support some of the claims and assertions made in the paper.

\textbf{Proficient:} \emph{Integrates references to support claims and assertions made in the paper.}

\textbf{Exemplary:} Integrates references to support claims and assertions made in the paper, effectively synthesizing different perspectives and research results from scholarly sources.

\begin{longtable}[]{@{}
  >{\raggedright\arraybackslash}p{(\columnwidth - 8\tabcolsep) * \real{0.2000}}
  >{\raggedright\arraybackslash}p{(\columnwidth - 8\tabcolsep) * \real{0.2000}}
  >{\raggedright\arraybackslash}p{(\columnwidth - 8\tabcolsep) * \real{0.2000}}
  >{\raggedright\arraybackslash}p{(\columnwidth - 8\tabcolsep) * \real{0.2000}}
  >{\raggedright\arraybackslash}p{(\columnwidth - 8\tabcolsep) * \real{0.2000}}@{}}
\toprule()
\begin{minipage}[b]{\linewidth}\raggedright
\textbf{TOTAL}
\end{minipage} & \begin{minipage}[b]{\linewidth}\raggedright
\textbf{0 = 0\% (F)}
\end{minipage} & \begin{minipage}[b]{\linewidth}\raggedright
\textbf{10 = 50\% (C)}
\end{minipage} & \begin{minipage}[b]{\linewidth}\raggedright
\textbf{15 = 75 (B)}
\end{minipage} & \begin{minipage}[b]{\linewidth}\raggedright
\textbf{20 = 100\% (A+)}
\end{minipage} \\
\midrule()
\endhead
\bottomrule()
\end{longtable}

\begin{center}\rule{0.5\linewidth}{0.5pt}\end{center}

\hypertarget{assignment-company-website-analysis}{%
\section*{Assignment: Company Website Analysis}\label{assignment-company-website-analysis}}
\addcontentsline{toc}{section}{Assignment: Company Website Analysis}

\begin{assessment}
Investigate the Human Resources or Faculty Development portion of a
company's website, a higher education institution or adult learning
facility, preferably one with which you are familiar. Focus on the
faculty or employee development part of the website. In this assignment,
you will apply the theory of teaching in/for/with depth by analyzing the
learning culture of an organization.

In a 4-5 page APA formatted paper, analyze the website by responding to
the following questions in your report:

\begin{enumerate}
\def\labelenumi{\arabic{enumi}.}
\tightlist
\item
  What can you infer about the company's learning culture?
\item
  From what is visible on the public website, would you say it is an
  authentic learning community? Why or why not? Discuss whether the
  website reflects aspects of one or more of the learning community
  models explored in previous lessons.
\item
  Do you see evidence that interconnectedness and integrity are valued?
  Explain.
\item
  What traits and skills seem to be valued in employees?
\item
  How does the company develop skills in its employees (e.g., workshops,
  seminars, mentoring)? Are the methods based on the principles of
  andragogy? (see Smith YouTube video). What specific adult learning
  strategies do you see reflected in the development/training
  opportunities for employees?
\end{enumerate}

Your paper should be 4-5 pages and should incorporate references to at
least five scholarly sources you have studied in this course, or other
scholarly sources you have identified.

The paper should include:

\begin{enumerate}
\def\labelenumi{\arabic{enumi}.}
\tightlist
\item
  Introduction
\item
  Analysis (responding to the prompts)
\item
  Conclusion
\item
  Reference List
\end{enumerate}
\end{assessment}

\hypertarget{company-website-analysis-rubric}{%
\subsection*{Company Website Analysis Rubric}\label{company-website-analysis-rubric}}
\addcontentsline{toc}{subsection}{Company Website Analysis Rubric}

See the following rubric that explains how your assignment will be evaluated. Also available as a \href{assets/assessment/Company-Website-Analysis-RUBRIC.pdf}{pdf}

\hypertarget{apa-formatting}{%
\subsubsection*{APA Formatting}\label{apa-formatting}}
\addcontentsline{toc}{subsubsection}{APA Formatting}

\textbf{Unsatisfactory:} Paper does not model language and conventions used in scholarly literature.
Writing is not well-organized. Several errors in grammar or composition. Sources
are not cited. APA citations are not appropriately formatted.

\textbf{Developing:} Paper partially models language and conventions used in scholarly literature.
Writing is somewhat well organized and includes some errors in grammar or
composition. Not all sources cited. APA citations are generally formatted
correctly, with several errors.

\textbf{Proficient:} \emph{Paper consistently models language and conventions used in scholarly
literature. Writing is well-organized and includes few (if any) errors in
grammar or composition. All resources are appropriately cited (including in-text
citations and bibliography information). Few (if any) errors in APA citations.}

\textbf{Exemplary:} Paper is an exemplar of language and conventions used in scholarly literature.
Writing is well-organized and free of errors in grammar or composition. All
resources are appropriately cited. No errors in APA format.

\hypertarget{developing-a-cohesive-and-logical-academic-argument-1}{%
\subsubsection*{DEVELOPING a COHESIVE and LOGICAL ACADEMIC ARGUMENT}\label{developing-a-cohesive-and-logical-academic-argument-1}}
\addcontentsline{toc}{subsubsection}{DEVELOPING a COHESIVE and LOGICAL ACADEMIC ARGUMENT}

\textbf{Unsatisfactory:} Does not make a focused, cohesive, or logical academic argument. Paper is
confusing, and is missing an introduction, body, or conclusion. Transitions
between sections and ideas are missing.

\textbf{Developing:} Makes an academic argument that is only partially focused, cohesive and logical.
Paper is generally organized, but is missing an introduction, body, or
conclusion. Transitions between sections and ideas are unclear.

\textbf{Proficient:} \emph{Makes a focused, cohesive, logical academic argument. Paper is effectively
organized and includes an introduction, body, and conclusion. Transitions
between sections and ideas are clear.}

\textbf{Exemplary:} Makes a focused, cohesive, logical and compelling academic argument. Paper is
effectively organized and includes an introduction, body, and conclusion.
Transitions between sections and ideas are clear and build on each other.

\hypertarget{analysis-of-learning-culture}{%
\subsubsection*{ANALYSIS of LEARNING CULTURE}\label{analysis-of-learning-culture}}
\addcontentsline{toc}{subsubsection}{ANALYSIS of LEARNING CULTURE}

\textbf{Unsatisfactory:} Does not include an analysis of the company learning culture, and no evaluation
of the authenticity of the learning community.

\textbf{Developing:} Includes a partial analysis of the company learning culture, including a limited
evaluation of the authenticity of the learning community.

\textbf{Proficient:} \emph{Includes a detailed analysis of the company learning culture, including an
evaluation of the authenticity of the learning community.}

\textbf{Exemplary:} Includes a detailed analysis of the company learning culture, including an
evaluation of the authenticity of the learning community. Includes a thoughtful
analysis, integrating scholarly literature to support analysis and furthering
scholarly thinking related to teacher identity.

\hypertarget{evaluation-of-interconnectedness-and-integrity}{%
\subsubsection*{EVALUATION of INTERCONNECTEDNESS and INTEGRITY}\label{evaluation-of-interconnectedness-and-integrity}}
\addcontentsline{toc}{subsubsection}{EVALUATION of INTERCONNECTEDNESS and INTEGRITY}

\textbf{Unsatisfactory:} Does not include an evaluation of evidence of interconnectedness and integrity
on the company website. Does not integrate scholarly sources in the evaluation.

\textbf{Developing:} Includes a partial evaluation of evidence of interconnectedness and integrity on
the company website. Evaluation includes only limited reference to scholarly
sources.

\textbf{Proficient:} \emph{Includes a detailed evaluation of evidence of interconnectedness and integrity
on the company website. Evaluation integrates scholarly sources.}

\textbf{Exemplary:} Includes a detailed evaluation of evidence of interconnectedness and integrity
on the company website. Includes recommendations for ways in which to integrate
interconnectedness and integrity into employee development.

\hypertarget{analysis-of-adult-learning-strategies}{%
\subsubsection*{ANALYSIS of ADULT LEARNING STRATEGIES}\label{analysis-of-adult-learning-strategies}}
\addcontentsline{toc}{subsubsection}{ANALYSIS of ADULT LEARNING STRATEGIES}

\textbf{Unsatisfactory:} Does not include a detailed analysis of valued skills and evidence of adult
learning theory in employee development. Does not integrate scholarly sources.

\textbf{Developing:} Includes a partial analysis of valued skills and evidence of adult learning
theory in employee development. Analysis integrates few, if any, scholarly
sources.

\textbf{Proficient:} \emph{Includes a detailed analysis of valued skills and evidence of adult learning
theory in employee development. Analysis integrates scholarly sources.}

\textbf{Exemplary:} Includes a detailed analysis of valued skills and evidence of adult learning
theory in employee development. Includes recommendations for ways in which to
integrate adult learning theory into employee development.

\hypertarget{scholarly-integration-1}{%
\subsubsection*{SCHOLARLY INTEGRATION}\label{scholarly-integration-1}}
\addcontentsline{toc}{subsubsection}{SCHOLARLY INTEGRATION}

\textbf{Unsatisfactory:} Does not integrate scholarly references to support claims and assertions made in
the paper.

\textbf{Developing:} Integrates scholarly references to support some of the claims and assertions
made in the paper.

\textbf{Proficient:} \emph{Integrates scholarly references to support claims and assertions made in the
paper.}

\textbf{Exemplary:} Integrates scholarly references to support claims and assertions made in the
paper, effectively synthesizing different perspectives and research results from
scholarly sources.

\begin{longtable}[]{@{}
  >{\raggedright\arraybackslash}p{(\columnwidth - 8\tabcolsep) * \real{0.2000}}
  >{\raggedright\arraybackslash}p{(\columnwidth - 8\tabcolsep) * \real{0.2000}}
  >{\raggedright\arraybackslash}p{(\columnwidth - 8\tabcolsep) * \real{0.2000}}
  >{\raggedright\arraybackslash}p{(\columnwidth - 8\tabcolsep) * \real{0.2000}}
  >{\raggedright\arraybackslash}p{(\columnwidth - 8\tabcolsep) * \real{0.2000}}@{}}
\toprule()
\begin{minipage}[b]{\linewidth}\raggedright
\textbf{TOTAL}
\end{minipage} & \begin{minipage}[b]{\linewidth}\raggedright
\textbf{0 = 0\% (F)}
\end{minipage} & \begin{minipage}[b]{\linewidth}\raggedright
\textbf{10 = 50\% (C)}
\end{minipage} & \begin{minipage}[b]{\linewidth}\raggedright
\textbf{15 = 75 (B)}
\end{minipage} & \begin{minipage}[b]{\linewidth}\raggedright
\textbf{20 = 100\% (A+)}
\end{minipage} \\
\midrule()
\endhead
\bottomrule()
\end{longtable}

\begin{center}\rule{0.5\linewidth}{0.5pt}\end{center}

\hypertarget{assignment-platform-paper}{%
\section*{Assignment: Platform Paper}\label{assignment-platform-paper}}
\addcontentsline{toc}{section}{Assignment: Platform Paper}

\begin{assessment}
For this assignment, you will write a contextualized Platform Paper in
which you discuss your ideal learning community and your role as
teacher/leader of that learning community. Select a context for your
paper (i.e.~facilitating in a FAR Centre in a specific country, teaching
adult learners, facilitating employee development workshops, etc.). Your
paper should be written and referenced in APA format and include
references to a minimum of 10 scholarly sources (this can include
literature you read in this course). You will write a draft of the
Platform Paper in Unit 8 and post for Peer Review. In Unit 9, you will
provide feedback to another learner on their paper. You will make
revisions based on the Peer Review and, in Unit 10, you will submit the
final Platform Paper. Peer reviewers will be assigned in advance.

\hypertarget{paper-outline}{%
\subsubsection{Paper Outline}\label{paper-outline}}

This paper will be 12-15 pages long, and should include: 1. Introduction
(1-2 pages) 2. Section 1: Ideal Learning Environment (5-7 pages) 3.
Section 2: Your Role as Teacher and Leader (5-7 pages) 4. Conclusion
(1-2 pages)

\hypertarget{paper-guidelines}{%
\subsubsection{Paper Guidelines}\label{paper-guidelines}}

\begin{itemize}
\tightlist
\item
  \textbf{Introduction}: Introduce the two sections in your paper,
  providing a brief description of the key points you will make in each
  section.
\item
  \textbf{Section 1}: In section one, you will describe your ideal
  education learning environment. This section should demonstrate your
  learning about authentic learning communities, incorporating scholarly
  sources and your own analysis to depict your ideal learning
  environment. Incorporate a discussion of the learning community
  environment, learning experiences, student learning outcomes, and
  personal beliefs about teaching and learning.
\item
  \textbf{Section 2}: In this section, describe your role as a teacher
  or leader within an authentic learning community. Incorporating
  scholarly literature, analyze your role as a facilitator/leader in
  planning learning experiences, facilitating student learning, and
  assessing student learning. Describe the actions, practices, and
  strategies you will engage in to achieve your vision of the learning
  community you described in section one.
\item
  \textbf{Conclusion}: Summarize the key points you made in each
  section.
\item
  \textbf{References}: Include a reference list with references to at
  least 10 scholarly sources.
\end{itemize}
\end{assessment}

\hypertarget{platform-paper-rubric}{%
\subsection*{Platform Paper Rubric}\label{platform-paper-rubric}}
\addcontentsline{toc}{subsection}{Platform Paper Rubric}

See the following rubric that explains how your assignment will be evaluated. Also available as a \href{assets/assessment/Platform-Paper-RUBRIC.pdf}{pdf}

\hypertarget{apawriting}{%
\subsubsection*{APA/WRITING}\label{apawriting}}
\addcontentsline{toc}{subsubsection}{APA/WRITING}

\textbf{Unsatisfactory:} Paper does not model language and conventions used in scholarly literature. Writing is not well-organized. Several errors in grammar or composition. Sources are not cited. APA citations are not appropriately formatted.

\textbf{Developing:} Paper partially models language and conventions used in scholarly literature. Writing is somewhat well organized and includes some errors in grammar or composition. Not all sources cited. APA citations are generally formatted correctly, with several errors.

\textbf{Proficient:} \emph{Paper consistently models language and conventions used in scholarly literature. Writing is well-organized and includes few (if any) errors in grammar or composition. All resources are appropriately cited (including in-text citations and bibliography information). Few (if any) errors in APA citations.}

\textbf{Exemplary:} Paper is an exemplar of language and conventions used in scholarly literature. Writing is well-organized and free of errors in grammar or composition. All resources are appropriately cited. No errors in APA format.

\hypertarget{developing-a-cohesive-and-logical-academic-argument-2}{%
\subsubsection*{DEVELOPING a COHESIVE and LOGICAL ACADEMIC ARGUMENT}\label{developing-a-cohesive-and-logical-academic-argument-2}}
\addcontentsline{toc}{subsubsection}{DEVELOPING a COHESIVE and LOGICAL ACADEMIC ARGUMENT}

\textbf{Unsatisfactory:} Does not make a focused, cohesive, or logical academic argument. Paper is confusing, and is missing an introduction, body, or conclusion. Transitions between sections and ideas are missing.

\textbf{Developing:} Makes an academic argument that is only partially focused, cohesive and logical. Paper is generally organized, but is missing an introduction, body, or conclusion. Transitions between sections and ideas are unclear.

\textbf{Proficient:} \emph{Makes a focused, cohesive, logical academic argument. Paper is effectively organized and includes an introduction, body, and conclusion. Transitions between sections and ideas are clear.}

\textbf{Exemplary:} Makes a focused, cohesive, logical and compelling academic argument. Paper is effectively organized and includes an introduction, body, and conclusion. Transitions between sections and ideas are clear, and build on each other.

\hypertarget{ideal-learning-environment}{%
\subsubsection*{IDEAL LEARNING ENVIRONMENT}\label{ideal-learning-environment}}
\addcontentsline{toc}{subsubsection}{IDEAL LEARNING ENVIRONMENT}

\textbf{Unsatisfactory:} Does not include a description of your ideal learning environment. Does not reference scholarly sources. Does note analyze key elements of an authentic learning community. Does not mention or describe the learning community environment, student learning outcomes, learning outcomes and personal beliefs about teaching and learning.

\textbf{Developing:} Includes a partial description of your ideal learning environment, referencing few scholarly sources and including a partial analysis of key elements of an authentic learning community. Mentions some elements, but does not fully describe the learning community environment, student learning outcomes, learning outcomes and personal beliefs about teaching and learning.

\textbf{Proficient:} \emph{Includes a detailed description of your ideal learning environment, referencing scholarly sources and analyzing key elements of an authentic learning community. Describes the learning community environment, student learning outcomes, learning outcomes and personal beliefs about teaching and learning.}

\textbf{Exemplary:} Includes a detailed description of your ideal learning environment, referencing scholarly sources and analyzing key elements of authentic learning communities. Provides a rationale for key elements of the learning community environment, student learning outcomes, learning outcomes and personal beliefs about teaching and learning. Advances scholarly thinking about authentic learning communities.

\hypertarget{your-role-as-teacher-and-leaders}{%
\subsubsection*{YOUR ROLE AS TEACHER AND LEADERS}\label{your-role-as-teacher-and-leaders}}
\addcontentsline{toc}{subsubsection}{YOUR ROLE AS TEACHER AND LEADERS}

\textbf{Unsatisfactory:} Does not include a description of your role as a teacher or leader within an authentic learning community, incorporating scholarly literature. Does not include an analysis of your role as a facilitator/leader in planning learning experiences, facilitating student learning, and assessing student learning. Does not include a description of the actions, practices, and strategies you will engage in to achieve your vision of the learning community you described in section one.

\textbf{Developing:} Includes a partial description of your role as a teacher or leader within an authentic learning community, incorporating scholarly literature. Describes but does not analyze your role as a facilitator/leader in planning learning experiences, facilitating student learning, and assessing student learning. Lists but does not describe the actions, practices, and strategies you will engage in to achieve your vision of the learning community you described in section one.

\textbf{Proficient:} \emph{Includes a detailed description of your role as a teacher or leader within an authentic learning community, incorporating scholarly literature. Includes a detailed analysis of your role as a facilitator/leader in planning learning experiences, facilitating student learning, and assessing student learning. Includes a detailed description of the actions, practices, and strategies you will engage in to achieve your vision of the learning community you described in section one.}

\textbf{Exemplary:} Includes a detailed analysis of your role as a teacher or leader within an authentic learning community, incorporating scholarly literature. Includes a detailed analysis of your role as a facilitator/leader in planning learning experiences, facilitating student learning, and assessing student learning. Includes a detailed description of the actions, practices, and strategies you will engage in to achieve your vision of the learning community you described in section one. Synthesizes scholarly thinking about the role of the teacher/leader.

\hypertarget{scholarly-integration-2}{%
\subsubsection*{SCHOLARLY INTEGRATION}\label{scholarly-integration-2}}
\addcontentsline{toc}{subsubsection}{SCHOLARLY INTEGRATION}

\textbf{Unsatisfactory:} Does not integrate many references to support the arguments made in the paper.

\textbf{Developing:} Integrates fewer than 10 scholarly sources to support arguments made in the paper.

\textbf{Proficient:} \emph{Integrates a minimum of 10 scholarly sources to support arguments made in each section of the paper.}

\textbf{Exemplary:} Integrates a minimum of 10 references to support the arguments made in each section, including several scholarly sources not included in course materials.

\begin{longtable}[]{@{}
  >{\raggedright\arraybackslash}p{(\columnwidth - 8\tabcolsep) * \real{0.2000}}
  >{\raggedright\arraybackslash}p{(\columnwidth - 8\tabcolsep) * \real{0.2000}}
  >{\raggedright\arraybackslash}p{(\columnwidth - 8\tabcolsep) * \real{0.2000}}
  >{\raggedright\arraybackslash}p{(\columnwidth - 8\tabcolsep) * \real{0.2000}}
  >{\raggedright\arraybackslash}p{(\columnwidth - 8\tabcolsep) * \real{0.2000}}@{}}
\toprule()
\begin{minipage}[b]{\linewidth}\raggedright
\textbf{TOTAL}
\end{minipage} & \begin{minipage}[b]{\linewidth}\raggedright
\textbf{0 = 0\% (F)}
\end{minipage} & \begin{minipage}[b]{\linewidth}\raggedright
\textbf{10 = 50\% (C)}
\end{minipage} & \begin{minipage}[b]{\linewidth}\raggedright
\textbf{15 = 75 (B)}
\end{minipage} & \begin{minipage}[b]{\linewidth}\raggedright
\textbf{20 = 100\% (A+)}
\end{minipage} \\
\midrule()
\endhead
\bottomrule()
\end{longtable}

\hypertarget{references}{%
\chapter*{References}\label{references}}
\addcontentsline{toc}{chapter}{References}

The following are key references used in this course. \textbf{\emph{Check with your course syllabus for required readings.}}

  \bibliography{book.bib}

\end{document}
